\subsection{映射的柱}

设 X 与 Y 是拓扑空间，f 是从 X 到 Y 的连续映射。用 U 表示空间 \(U_{1} = X \times I\) 与空间 \(U_{2} = Y\) 的拓扑和空间，并经由典范单射把 \(U_{1}\) 与 \(U_{2}\) 等同于 U 的子空间。设 R 是 U 上最细的等价关系，使得 \((x,1)\)（\(U_{1}\) 的点）与 \(f(x)\)（\(U_{2}\) 的点）等价（对每个 \(x \in X\)）。

定义 7 —— 我们把商拓扑空间 U/R 称为映射 f 的柱，记作 \(\mathrm{Cyl}(f)\)。

用 \(\alpha_{f}\colon\mathrm{X}\times\mathbf{I}\to\mathrm{Cyl}(f)\) 与 \(\beta_{f}\colon\mathrm{Y}\to\mathrm{Cyl}(f)\) 表示从 U 到 U/R 上的典范满射在 \(\mathrm{U}_{1}\) 与 \(\mathrm{U}_{2}\) 上的限制。映射 \(\alpha_{f}\) 是一个同伦，其终点为 \(\beta_{f}\circ f\)。

命题 4（柱的泛性质）

设 Z 是拓扑空间，\(\beta: \mathrm{Y} \to \mathrm{Z}\) 是连续映射，\(\alpha: \mathrm{X} \times \mathbf{I} \to \mathrm{Z}\) 是终点为 \(\beta \circ f\) 的同伦。存在唯一的连续映射 \(\varphi\)（从 \(\operatorname{Cyl}(f)\) 到 Z），使得 \(\alpha = \varphi \circ \alpha_f\) 且 \(\beta = \varphi \circ \beta_f\)。

映射 \(\psi\)（从 U 到 Z，与 \(\alpha\) 在 \(\times\) 上、与 \(\beta\) 在 Y 上一致）是连续的。我们有 \(\alpha(x,1)=\beta(f(x))\)（对所有 \(x\in X\)），因此 \(\psi\) 过渡到商，给出连续映射 \(\varphi\)（从 \(\mathrm{Cyl}(f)\) 到 Z），使得 \(\alpha=\varphi\circ\alpha_{f}\) 且 \(\beta=\varphi\circ\beta_{f}\)。由于 \(\alpha_{f}\) 与 \(\beta_{f}\) 的像覆盖 \(\mathrm{Cyl}(f)\)，\(\varphi\) 是满足这些关系式的从 \(\mathrm{Cyl}(f)\) 到 Z 的唯一映射。

例 1 —— 设 X' 与 Y' 是拓扑空间，\(f': \mathrm{X}' \to \mathrm{Y}'\) 是连续映射。设 \(u: \mathrm{X} \to \mathrm{X}'\) 与 \(v: \mathrm{Y} \to \mathrm{Y}'\) 是满足 \(f' \circ u = v \circ f\) 的连续映射。存在唯一的连续映射 \(w: \operatorname{Cyl}(f) \to \operatorname{Cyl}(f')\)，使得 \(w \circ \alpha_f = \alpha_{f'} \circ (u \times \operatorname{Id}_\mathbf{I})\) 且 \(w \circ \beta_f = \beta_{f'} \circ v\)。

把命题 4 应用于 Z=Y 与 \(\beta = \operatorname{Id}_{Y}\)，存在唯一的连续映射 \(\gamma_{f}\colon\operatorname{Cyl}(f)\to\operatorname{Y}\)，使得有 \(\gamma_{f}(\alpha_{f}(x,s))=f(x)\) 与 \(\gamma_{f}(\beta_{f}(y))=y\)（对 \(x\in X\)、\(s\in I\) 及 \(y\in Y\)）。

命题 5 —— 映射 \(\alpha_{f}\) 诱导 X × {[}0,1{[} 到 \(\mathrm{Cyl}(f)\) 的一个开子集上的同胚。映射 \(\beta_{f}\) 定义 Y 到该开集的余集上的同胚。映射 \(\beta_{f} \circ \gamma_{f}\) 是 \(\mathrm{Cyl}(f)\) 到 \(\beta_{f}(\mathrm{Y})\) 上的连续收缩映射。

集合 X × {[}0,1{[} 是 U 的开子集，对关系 R 是饱和的，且 R 在 X × {[}0,1{[} 中诱导的等价关系是相等关系。由此即得第一个断言（TG, I, p. 23, 推论 10）。

\(\alpha_{f}(\mathrm{X}\times[0,1])\) 的余集是 \(\beta_{f}(\mathrm{Y})\)。由于 \(\gamma_{f}\circ\beta_{f}=\mathrm{Id}_{\mathrm{Y}}\)，\(\beta_{f}\) 定义 Y 到 \(\beta_{f}(\mathrm{Y})\) 上的同胚，且 \(\beta_{f}\circ\gamma_{f}\) 是 \(\beta_{f}(\mathrm{Y})\) 到 \(\mathrm{Cyl}(f)\) 中的典范单射的连续收缩映射。

\(\beta_{f}(\mathrm{Y})\)（\(\operatorname{Cyl}(f)\) 的闭子空间）称为 \(f\) 的柱的柱底。映射 \(\beta_{f} \circ \gamma_{f}\) 称为 \(\operatorname{Cyl}(f)\) 到其柱底上的典范收缩映射。

考虑由下式定义的映射 \(\sigma_{1}\colon\mathrm{U}_{1}\times\mathbf{I}\to\mathrm{Cyl}(f)\) 与 \(\sigma_{2}\colon\mathrm{U}_{2}\times\mathbf{I}\to\mathrm{Cyl}(f)\)：

\[
\begin{array}{c c} \sigma_ {1} ((x, s), t) = \alpha_ {f} (x, (1 - t) s + t) & \text {for} (x, s) \in \mathrm{X} \times \mathbf {I} \text {and} t \in \mathbf {I} \\ \sigma_ {2} (y, t) = \beta_ {f} (y) & \text {for} y \in \mathrm{Y} \text {and} t \in \mathbf {I}. \end{array}
\]

它们是连续的。对 \(x \in X\) 及 \(t \in I\)，我们有

\[
\sigma_ {1} ((x, 1), t) = \alpha_ {f} (x, 1) = \beta_ {f} (f (x)) = \sigma_ {2} (f (x), t).
\]

因此存在唯一的映射 \(\sigma_{f}\)（从 \(\mathrm{Cyl}(f)\times\mathbf{I}\) 到 \(\mathrm{Cyl}(f)\)），使得 \(\sigma_{f}\circ(\alpha_{f}\times\mathrm{Id}_{\mathbf{I}})=\sigma_{1}\) 且 \(\sigma_{f}\circ(\beta_{f}\times\mathrm{Id}_{\mathbf{I}})=\sigma_{2}\)。映射 \(\sigma_{f}\) 是连续的（I, p. 19, 命题 10）。

命题 6 —— 映射 \(\sigma_{f}\colon\mathrm{Cyl}(f)\times\mathbf{I}\to\mathrm{Cyl}(f)\) 是 \(\mathrm{Cyl}(f)\) 到其柱底上的一个强收缩。其终点是 \(\mathrm{Cyl}(f)\) 到其柱底上的典范收缩映射。

映射 \(\sigma_f\) 是一个同伦。关系式 \(\sigma_f(\alpha_f(x,s),0) = \alpha_f(x,s)\) 与 \(\sigma_f(\beta_f(y),0) = \beta_f(y)\) 表明 \(\sigma_f\) 的起点是 \(\mathrm{Cyl}(f)\) 的恒同映射。用 \(r_f\) 表示 \(\sigma_f\) 的终点。关系式

\[
\sigma_ {f} \left(\alpha_ {f} (x, s), 1\right) = \alpha_ {f} (x, 1) = \beta_ {f} (f (x)) = \left(\beta_ {f} \circ \gamma_ {f}\right) \left(\alpha_ {f} (x, s)\right)
\]

与 \(\sigma_{f}(\beta_{f}(y),1) = \beta_{f}(y) = (\beta_{f}\circ \gamma_{f})(\beta_{f}(y))\) 表明 \(r_f\) 是典范收缩映射 \(\beta_{f}\circ \gamma_{f}\)（\(\mathrm{Cyl}(f)\) 到其柱底上的）。

对 \((x,s)\in \mathbf{X}\times \mathbf{I}\) 及 \(t\in \mathbf{I}\)，我们有

\[
r _ {f} \left(\sigma_ {f} \left(\alpha_ {f} (x, s), t\right)\right) = r _ {f} \left(\alpha_ {f} (x, s (1 - t) + t)\right) = \beta_ {f} (f (x)) = r _ {f} \left(\alpha_ {f} (x, s)\right).
\]

对 \(y \in Y\) 及 \(t \in I\)，有 \(r_{f}(\sigma_{f}(\beta_{f}(y), t)) = r_{f}(\beta_{f}(y))\)。于是，\(\sigma_{f}\) 是 \(\operatorname{Cyl}(f)\) 到其柱底上的强收缩。

映射 \(\sigma_{f}\) 称为 \(f\) 的柱到其柱底上的典范收缩。

注 —— 1) 设 A 是 X × I 的子集，A₁ 是使 (x,1) ∈ A 的点 x ∈ X 所成的集合。那么我们有 \(\bar{\alpha}_{f}^{-1}(\alpha_{f}(A)) = A \cup f^{-1}(f(A_{1})) \times \{1\}\) 与 \(\beta_{f}^{-1}(\alpha_{f}(A)) = f(A_{1})\)。因此，若 f 是闭（相应地，开）的，则映射 \(\alpha_{f}\) 亦然。

\begin{enumerate}
\def\labelenumi{\arabic{enumi})}
\setcounter{enumi}{1}
\tightlist
\item
  设映射 \(f\) 是逆紧的。设 \(P\) 是 \(\operatorname{Cyl}(f)\) 的一点。若 \(P = \alpha_f(x,t)\)，其中 \(0 \leqslant t < 1\) 且 \(x \in X\)，则我们有 \(\overline{\alpha}_f^1(P) = (x,t)\)。在相反情形，存在 \(y \in Y\) 使得 \(P = \beta_f(y)\)
\end{enumerate}

且 \(\bar{\alpha}_{f}^{1}(P) = \bar{f}^{1}(y) \times \{1\}\)。这证明映射 \(\alpha_{f}\) 的纤维是拟紧的。由于 \(\alpha_{f}\) 是闭映射，故它是逆紧的（TG, I, p. 75, 定理 1）。

由命题 5 及 TG, I, p. 72, 命题 2，映射 \(\beta_f\) 本身也是逆紧的。因此，若 f 逆紧，则从 U 到 \(\mathrm{Cyl}(f)\) 上的典范满射是逆紧的，从而特别地是普遍严格的（I, p. 20, 推论）。

\begin{enumerate}
\def\labelenumi{\arabic{enumi})}
\setcounter{enumi}{2}
\tightlist
\item
  若空间 X 与 Y 是分离的，则 f 的柱也是分离的。事实上，设 z 与 \(z'\) 是 \(\mathrm{Cyl}(f)\) 的两个不同点，我们来证明它们有互不相交的邻域。我们区分三种情形。
\end{enumerate}

- 存在 \((x,t)\) 与 \((x',t') \in \mathrm{X} \times [0,1[\)，使得 \(z = \alpha_f(x,t)\) 且 \(z' = \alpha_f(x',t')\)。

在此情形，断言由如下事实得出：空间 \(X \times [0,1]\) 是分离的（TG, I, p. 54, 命题 7），且映射 \(\alpha_{f}\) 诱导该空间到 \(\operatorname{Cyl}(f)\) 的一个开子空间上的同胚。

- 存在 \((x,t)\in \mathrm{X}\times [0,1[\) 使得 \(z = \alpha_{f}(x,t)\)，且存在 \(y^{\prime}\in \mathrm{Y}\) 使得 \(z^{\prime} = \beta_{f}(y^{\prime})\)。

此时，\(\alpha_{f}(\mathrm{X}\times [0,\frac{t + 1}{2}])\) 与 \(\operatorname {Cyl}(f)\mathbf{-}\alpha_{f}(\mathrm{X}\times [0,\frac{t + 1}{2}])\) 是 \(z\) 与 \(z^{\prime}\) 在 \(\operatorname {Cyl}(f)\) 中的互不相交的开邻域。

- 存在 \(y\) 与 \(y'\in Y\)，使得 \(z=\beta_f(y)\) 且 \(z'=\beta_f(y')\)。

在此情形，\(y\neq y'\)；由于 Y 是分离的，存在 y 在 Y 中的开邻域 V 以及开邻域 \(V'\)（\(y'\) 在 Y 中的），使得 \(V\cap V'=\varnothing\)。于是，\((\beta_f\circ\gamma_f)^{-1}(V)\) 与 \((\beta_f\circ\gamma_f)^{-1}(V')\) 是 \(\operatorname{Cyl}(f)\) 的互不相交的开子集，分别包含 \(\beta_f(y)\) 与 \(\beta_f(y')\)。

